% Technical report (LaTeX version of docs/technical_report.md).
% Build: pdflatex technical_report.tex   (or: latexmk -pdf technical_report.tex)
\documentclass[10pt]{article}
\usepackage[a4paper,margin=2cm]{geometry}
\usepackage[T1]{fontenc}
\usepackage{lmodern}
\usepackage{amsmath}
\usepackage{booktabs}
\usepackage{enumitem}
\usepackage[hidelinks]{hyperref}
\usepackage{titlesec}

\setlist{nosep,leftmargin=1.4em}
\titlespacing*{\section}{0pt}{1.2ex plus .5ex}{.6ex}
\setlength{\parindent}{0pt}
\setlength{\parskip}{.6ex}

\newcommand{\pct}{\,\%}
\newcommand{\para}[1]{\textbf{#1}}

\title{\vspace{-1.5em}Learning drywall panel placement from UMI-style demonstrations\\
and simulation-based offline RL}
\author{Bruce Peters}
\date{\vspace{-1em}\small Technical report (draft). Repository: this repo; full log of decisions and failures:
\texttt{LABNOTEBOOK.md}.}

\begin{document}
\maketitle

\section{System architecture and design decisions}

\para{Simulation (Part A).} The simulation is a MuJoCo 4 work cell with a Franka Panda (from Menagerie, with
no hand) carrying a $24 \times 18$\,cm suction plate, a $0.6 \times 0.4 \times 0.01$\,m panel (1--3\,kg), a
rack, and a wall with a target region and a support ledge. Suction is a weld constraint that only engages when
the plate is within 3\,cm of the panel face, centred and within $30^\circ$ of parallel, and it locks in at the
\emph{current} relative pose (no snapping the panel into place). Once a released panel is resting on the
ledge within 2\,mm of the wall, a second weld holds it there as an external fixture.

The policy acts at 10\,Hz with a 7-D action: base-frame end-effector deltas (at most 2\,cm and $3^\circ$ per
step) plus a suction bit. The deltas are integrated into a commanded pose, clipped to a workspace box, and
tracked by damped-least-squares IK feeding position actuators. On real hardware the same commanded pose would
just be streamed to the vendor's Cartesian impedance controller.

Observations are split into \emph{deployable} ones (joint states, EE pose, suction state, and panel and target
pose \emph{estimates}) and \emph{privileged} ones (true poses, velocities, contact forces). Every reward term is
logged separately: position and orientation error, collisions, wall force above 50\,N, action magnitude and
jerk, a drop ($-50$) and success ($+100$). Success means the panel is within 20\,mm / $5^\circ$ of the target,
against the wall, released, and still for 5 steps.

\para{Two task distributions.} In the easy task the rack and wall are fixed, with $\pm 5$\,cm / $\pm 10^\circ$
randomization plus mass and friction. In the hard task the wall can rotate anywhere within $\pm 150^\circ$
around the robot and the rack can sit anywhere on a 0.6--0.7\,m ring (even behind the robot). I picked those
ranges from an IK reachability sweep, which showed the arm can only face the wall at panel heights of about 0.55--0.75\,m,
and a $+15\ldots+75^\circ$ band of panel yaw is unreachable at the wrist. The hard task needed two additions to
work at all, a posture-aware IK (the null-space posture follows the EE azimuth, otherwise joint 6 hits its
limit at rotated walls) and a 500-step limit.

\para{Demonstrations (Part B).} A scripted waypoint controller stands in for the human demonstrator, and it
only uses the action interface and deployable observations. The hard version works in the wall's frame,
swings around the base, and has two modes with per-episode setpoint variance: \texttt{demo} (moderate, mostly
successful) and \texttt{explore} (large variance and smooth noise, so a lot of failures).

Episodes are recorded as a UMI dataset in UMI's own zarr replay-buffer format (JpegXL images,
\texttt{episode\_ends}). Each one has 60\,Hz absolute EE poses in a world/SLAM frame, a gripper (suction)
state, a $155^\circ$ fisheye wrist camera (with mirrors), an external camera, the panel pose in the target
frame, episode boundaries and success labels. No actions or rewards are stored.

\para{UMI $\to$ robot conversion.} Turning the UMI data into robot training data takes five steps:
\begin{enumerate}
  \item World $\to$ robot-base frame (tested with an arbitrary rotated and translated map frame).
  \item Timestamp resampling to the 10\,Hz control rate (linear position, slerp rotation, zero-order-hold
    gripper).
  \item Absolute poses $\to$ base-frame deltas scaled to the action limits.
  \item Gripper width $\to$ suction bit.
  \item A workspace check.
\end{enumerate}

The learning state is \para{37-D} (EE pose, suction, panel and target poses, panel-in-target pose), and
\para{46-D} for the hard task (which adds the panel pose in the EE frame). All rotations use the continuous 6D
representation from Zhou et al.\ (2019), because the tool-down EE quaternion sits at $w \approx 0$ and
flipped sign 10 times in a single episode under a $w \ge 0$ convention.

To validate the conversion, I rebuilt states from the zarr and they match the simulator to $10^{-7}$.
Replaying the derived actions open-loop also reproduces the recorded outcomes: easy 73\pct{} vs 74\pct, and
hard 99.6\pct{} vs 99.8\pct{} with 99.4\pct{} outcome agreement.

This does make some big assumptions: no SLAM drift, perfect gripper sensing, the scripted controller as the
human, and a target pose straight from the simulator (standing in for fiducial or BIM registration). Panel and
target poses are perfect estimates, and no camera-based perception is implemented.

\section{Training methodology}

\para{Policy.} Every learned policy uses the same network, an MLP ($3 \times 256$ with LayerNorm) on
normalized states. It outputs 6 tanh motion dims and 1 suction logit as \para{action chunks} of 8 steps, of
which 4 are executed before re-planning. The loss is MSE + BCE (the negative log-likelihood of a
fixed-variance Gaussian $\times$ Bernoulli). Chunking was the single biggest change on the easy task, where
single-step BC only reached 16\pct{} validation success and chunking got that to 84\pct{} (it deals with phase
ambiguity and compounding error).

\para{Behaviour cloning} is trained on successful episodes only (``filtered BC'') for 30k steps.

\para{Offline RL.} I implemented Implicit Q-Learning (Kostrikov et al.\ 2021) \para{from scratch}, with each
loss documented against the paper's equations:
\begin{itemize}
  \item expectile value regression ($\tau = 0.7$);
  \item TD targets bootstrapped from $V(s')$, never from $\max_{a'} Q$;
  \item twin target critics;
  \item advantage-weighted regression for the actor with $\exp(\beta A)$, clipped at 100,
    $\beta \in \{3, 10\}$.
\end{itemize}
The actor \emph{is} the BC network, so BC and IQL only differ in the per-sample weight. Training only reads
the fixed dataset (\texttt{env.step} is never called for learning).

The \para{hard-task IQL fixes} came out of a diagnosis (Section~\ref{sec:failures}): a reward scale of 0.01
instead of 0.1, 50/50 demo/sim batches, release targets weighted $\times 10$ in the suction loss, and 1M
steps. Optionally, the release threshold $p(\text{attach})$ is tuned, on validation seeds only.

\para{Offline datasets.} The easy dataset is 1,000 episodes from scripted plus noisy-BC behaviour
(57\pct{} success), and the hard one is 2,000 \texttt{explore} episodes (64\pct{} success, 10\pct{} drops,
collisions in 41\pct). Both are stored as D4RL-style transitions with every reward term kept separately. When
UMI and sim data are combined, both get relabelled to only the reward terms UMI data can support, since UMI
has no collision, force or drop sensing.

\para{Protocol.} Data, validation and test seeds are all disjoint (the data scripts refuse test seeds).
Validation runs 50 episodes every 5k (BC) / 20k (IQL) steps and is only a measurement. Snapshots are saved for
inspection but never used to pick a checkpoint, so the reported model is always the last step. Testing is 100
episodes per policy with 95\pct{} Wilson intervals and an automatic failure-mode classifier. The hard task has
two test sets, one \emph{inside} the UMI demos' region and one \emph{outside} it.

\section{Experimental results}

\para{Easy task} (test seeds 10000--10099):

\begin{center}
\begin{tabular}{llll}
\toprule
policy & data & success (95\pct{} CI) & collisions \\
\midrule
scripted controller & --- & 100\pct{} (96--100) & 0\pct \\
BC, single-step & UMI demos & 23\pct{} (16--32) & 2\pct \\
\textbf{BC, chunked} & UMI demos & \textbf{80\pct{} (71--87)} & 9\pct \\
BC control & UMI + sim successes & 69\pct{} (59--77) & 2\pct \\
IQL $\beta 3$ / $\beta 10$ & sim & 58\pct{} / 61\pct & 22 / 14\pct \\
\textbf{IQL $\beta 3$} / $\beta 10$ & UMI + sim & \textbf{71\pct} / 68\pct & 8 / 12\pct \\
IQL $\beta 3$, release threshold tuned (0.65) & UMI + sim & \textbf{81\pct{} (72--87)} & 8\pct \\
\bottomrule
\end{tabular}
\end{center}

\para{Hard task} (in = seeds 30000--30099 inside the demos' region; out = 31000--31099 outside):

\begin{center}
\begin{tabular}{llll}
\toprule
policy & data & in & out \\
\midrule
scripted controller (demo mode) & --- & 99\pct{} (95--100) & 92\pct{} (85--96) \\
BC, chunked & UMI demos (demo region) & 6\pct{} (3--12) & \textbf{0\pct{} (0--4)} \\
\textbf{BC control} & UMI + sim successes & \textbf{55\pct{} (45--64)} & 20\pct{} (13--29) \\
IQL $\beta 10$, success $100 \to 1000$ & sim & 2\pct{} $\to$ 16\pct & 10\pct{} $\to$ 19\pct \\
IQL $\beta 3$, success $100 \to 1000$ & sim & 21\pct{} $\to$ 8\pct & 11\pct{} $\to$ 14\pct \\
IQL $\beta 10$, success 1000 & UMI + sim & 39\pct{} (30--49) & 22\pct{} (15--31) \\
IQL $\beta 10$, same, 1M steps (ablation) & UMI + sim & 33\pct{} (25--43) & 17\pct{} (11--26) \\
\textbf{IQL fixed $\beta 3$} / $\beta 1$ / $\beta 10$ & UMI + sim, balanced
  & \textbf{65\pct{} (55--74)} / 50 / 57\pct & 26 / \textbf{32\pct} / 19\pct \\
\bottomrule
\end{tabular}
\end{center}

Charts are in \texttt{media/*/comparison*.png} and \texttt{media/*/training\_*.png}, and per-episode results
are in \texttt{results/}.

\para{Findings.}
\begin{enumerate}[itemsep=.4ex]
  \item \para{Simulation data gives coverage that demonstrations don't.} BC on demos from a limited region
    scores 0\pct{} outside it (91 of 100 episodes never even grasp the panel). Every policy trained with
    simulated data generalizes outside the demo region, 20--32\pct{} for the best ones.
  \item \para{Untuned offline RL did not beat behaviour cloning, but tuned, it matched or edged past it.}
    Before tuning, IQL got 71\pct{} on the easy task vs 69\pct{} for the BC control and 80\pct{} for BC on
    demos, and on the hard task 39\pct{} / 22\pct{} vs 55\pct{} / 20\pct. After the fixes, hard IQL reached
    65\pct{} / 26\pct. Paired on the same test seeds, that is $+36$ / $-10$ episodes vs the old IQL
    ($p < 0.001$) and $+31$ / $-21$ vs the BC control ($p = 0.21$). On the easy task, a validation-tuned
    release threshold lifted IQL to 81\pct{} ($+10$ / $-0$, $p = 0.002$). One caveat: at the new reward scale
    the $\beta 3$ weights are nearly uniform (mean 1.01), so a lot of the gain might come from the balanced,
    release-weighted data rather than from the advantage weighting itself.
  \item \para{Combining demos with sim data helps IQL} in both tasks (easy 71 vs 58--61\pct; hard 39 vs
    16\pct{} in-region). For plain BC it depends on the task. On the easy task adding the noisier sim
    successes hurt ($80 \to 69$\pct), while on the hard task it was essential ($6 \to 55$\pct).
  \item \para{Reward balance matters, but it is not the whole story.} With success worth $+100$, the
    discounted success ($\approx 22$ from 150 steps out) was only worth about 70 steps of moderate error, so
    IQL learned to just keep holding the panel. Raising success $\times 10$ improved $\beta 10$
    ($2 \to 16$\pct{} in-region, with disjoint intervals) but not $\beta 3$ ($21 \to 8$\pct).
\end{enumerate}

\section{Failure analysis}\label{sec:failures}

\para{Where the policies fail.} Nearly every failure of the learned policies happens in the last few
centimetres or at the suction decision. Easy BC either stalled about 9\,mm above the panel without grasping (a regression fixed point, since the tiny
demonstrated motions near the goal average out to zero) or released 7\,mm short of the wall. Easy IQL fixed the
grasp (never-grasped went from 16--27 to 3--9) but then released off the wall. On the hard task, BC on demos
released mid-air (50 drops) or toggled suction at the target. The sim-trained policies were dominated by
``held at end'', where the panel gets carried but never placed. When I measured it, that turned out to be
mostly an approach failure: IQL's held episodes end a median 10--20\,cm from the target, while the BC
control's end 3\,cm away. It happens in BC too (which has no reward), so reward hacking can't be the whole explanation.

\para{Diagnosing hard IQL} is what led to the fixes. With success 1000 at reward $\times 0.1$ the weights
were saturated, with 11--17\pct{} of samples hitting the $\exp(\beta A)$ clip. Release was rare, only 1 in
144--224 suction-on states. And the runs were under-trained, with validation still rising at 200k steps.
Training longer on its own did not help ($39 \to 33$\pct), but the combination of fixes did (65\pct).

\para{Critic convergence.} At 200k steps the values are still drifting and validation swings
$\pm$15--20\pct{} between checkpoints. At 1M steps the values are nearly flat, but validation still swings
about $\pm 10$\pct.

\para{Collisions.} IQL policies collide more than the BC control trained on the same data (hard: 36--70\pct{}
vs 19--34\pct{} of episodes). They imitate the collision-heavy exploration data with only partial advantage
reweighting. BC on the demos alone collides the most (65--84\pct), because it is mostly outside its training
distribution.

\para{Other limitations:} a rigid panel and ideal suction, perfect pose estimates, one training seed per
configuration, and about 0.2\pct{} of hard sim episodes that hit a MuJoCo instability and were reset
mid-episode (detected after the fact and documented).

\section{Proposed deployment approach}

This is summarized from \texttt{docs/deployment\_proposal.md}.

\para{Hybrid architecture.} The learned policy (10\,Hz, 8-step chunks, 4 executed) handles gross motion,
carrying and the approach, streamed as Cartesian pose targets to a 1\,kHz impedance controller. The final step
gets verified instead of learned: the final approach becomes a force-controlled push to a set point, and the
suction release becomes a \emph{request} that is only granted once contact and support are confirmed. My
results show this is exactly where every learned policy fails, while explicit logic succeeds 92--100\pct{} of
the time.

\para{Sim-to-real risks.} The main ones (with how I would mitigate them) are perception (fiducials,
registration, tracking through the grasp), contact dynamics (force control, system identification), panel
flex and suction slip (vacuum monitoring, a multi-cup frame or the second arm), and site localization
(registration against BIM, a reachability check before acting).

\para{Failure handling.} If the panel slips the robot does a safe set-down, if the wall doesn't match it
re-registers or escalates, if contact comes early it backs off and re-aligns, and if the policy is uncertain
(ensemble disagreement, a low critic value) it pauses and hands over.

\para{Field adaptation with 10--20 UMI demos.} I would co-train on the site demos plus prior data, build a
digital twin of the site for evaluation, and then go up a staged validation ladder: open-loop replay $\to$
twin $\to$ dry run $\to$ dummy panel $\to$ supervised real panels. Each stage has explicit pass gates (for
example $\ge 95$\pct{} success with a Wilson lower bound $\ge 85$\pct, and zero safety violations).

\para{Bottom line.} A small number of human demonstrations teaches the motion, but simulation is what extends
it to layouts the demonstrations never covered. In this prototype offline RL only added success after careful
diagnosis and tuning, and not significantly beyond BC. The safety-critical contact and release should not be
left to the learned policy.

\end{document}
